\chapter{Factibilidad}
\section{Factibilidad técnica}
\subsection*{Factibilidad técnica - Hardware}

\subsection*{Factibilidad técnica - Software}
.

\section{Factibilidad Operativa}

\section{Factibilidad Económica}


\subsection{Estimación de ingresos o beneficios}

El sistema generará ahorros significativos al optimizar procesos operativos:

\begin{enumerate}
    \item \textbf{Ahorro por Reducción de Horas Extra en Producción.}  
    
Actualmente, debido a la falta de información precisa sobre los volúmenes diarios de producción requeridos,
el personal de producción debe extender su jornada laboral para realizar ajustes de última hora.
En promedio, esta situación genera una hora extraordinaria diaria para 3 panaderos y 1 cocedor de pan,
quienes cumplen una jornada regular de 40 horas semanales.
Para efectos del análisis económico, se considera que esta extensión de jornada ocurre en promedio
3 días a la semana, lo que equivale a 12 días al mes, según antecedentes operativos históricos de la Panadería Oriente.

% horas extra mensuales
\[
\underbrace{1}_{\text{h/día}} \times
\underbrace{4}_{\text{trabajadores}} \times
\underbrace{12}_{\text{días/mes}}
= 48 \text{ h/mes}
\]

% valor hora extra (40 horas semanales)
\[
\text{Valor hora extra} =
\left(
\frac{750.000}{30}
\times
\frac{28}{160}
\right)
\times 1{,}5
=
\text{\$}6.563/\text{h}
\]

% ahorro económico
\[
\begin{aligned}
\text{Ahorro mensual} &= 48 \times \text{\$}6.563 = \text{\$}315.000 \\[4pt]
\Rightarrow\ \text{Ahorro anual} &= \text{\$}315.000 \times 12 = \text{\$}3.780.000
\end{aligned}
\]

\textit{El cálculo del valor de la hora extraordinaria se realiza conforme al Artículo 32 del \textcite{codigotrabajo} chileno. }

    

    
\item \textbf{Reducción de Merma en Repartos.}

A partir del análisis del promedio de los últimos 6 meses, se identificó que la merma total de la panadería se compone de pérdidas por repartos, barrido y ración de pan. De estas, únicamente la merma asociada a repartos es susceptible de ser reducida mediante la implementación del sistema propuesto.

La merma acumulada en repartos asciende a \$3.200.688. Considerando un escenario conservador, se estima que el sistema podría contribuir a una reducción del 20\% de dicha merma:

\[
\begin{aligned}
\text{Ahorro mensual} &= \text{\$}3.200.688 \times 0.20 = \text{\$}960.206 \\[4pt]
\Rightarrow\ \text{Ahorro anual} &= \text{\$}960.206 \times 12 = \text{\$}7.681.656
\end{aligned}
\]


    
    \item \textbf{Ahorro en Supervisor de Reparto.}
    De 8 visitas mensuales se espera que el supervisor solo haga la mitad de ellas a \$20.000 cada una:
    \[
    4 \times \text{\$}20.000 = \text{\$}80.000\ \text{por mes}
    \qquad\Rightarrow\qquad
    \text{Ahorro anual} = \text{\$}80.000 \times 12 = \text{\$}960.000
    \]
\end{enumerate}

\noindent\textbf{Resumen de ahorros anuales:}
\[
\text{\$}3.780.000 + \text{\$}7.681.656 + \text{\$}960.000  
= \boxed{\text{\$}12.421.656}
\]


%ANALISIS DEL VAN (FORMULA + TABLA )
\subsection{Análisis del Valor Actual Neto (VAN)}

Para evaluar la viabilidad económica, se calcula el VAN considerando los siguientes parámetros:

\begin{itemize}
    \item Vida útil del sistema: \textbf{5 años.}
    \item Tasa de descuento: \textbf{12\%.}
\end{itemize}

La vida útil del sistema se estima en función de la obsolescencia tecnológica, utilizando tecnologías actuales con amplia comunidad de soporte. 
También se considera la naturaleza del negocio y el plan del cliente como expectativa mínima de duración del software.

\vspace{0.8em}
\noindent Para calcular el VAN se utiliza la siguiente ecuación \textcite{investopedia}:

\[
VAN = \sum_{t=1}^{n} \frac{V_t}{(1+k)^t} - I_o
\]

\noindent Donde:
\begin{itemize}
    \item $V_t$: Flujo de caja neto en el periodo $t$.
    \item $I_o$: Inversión inicial.
    \item $n$: Número de periodos.
    \item $k$: Tasa de descuento.
\end{itemize}

\vspace{1em}
\begin{table}[H]
\centering
\renewcommand{\arraystretch}{1.3}
\setlength{\tabcolsep}{6pt}
\begin{tabular}{|c|c|c|c|c|c|c|}
\hline
\rowcolor[HTML]{EAEAEA}
\textbf{VAN} & \textbf{Año 0} & \textbf{Año 1} & \textbf{Año 2} & \textbf{Año 3} & \textbf{Año 4} & \textbf{Año 5} \\ \hline
\textbf{Ingresos} &  & \$12.421.656 & \$12.421.656 & \$12.421.656 & \$12.421.656 & \$12.421.656 \\ \hline
\textbf{Gastos} &  & -\$835.996 & -\$835.996 & -\$835.996 & -\$835.996 & -\$835.996 \\ \hline
\textbf{Inversión} & \$7.475.196 &  &  &  &  &  \\ \hline
\textbf{Flujo} & -\$7.475.196 & \$11.585.660 & \$11.585.660 & \$11.585.660 & \$11.585.660 & \$11.585.660 \\ \hline
\multicolumn{6}{|r|}{\cellcolor[HTML]{EAEAEA}\textbf{VAN =}} & \cellcolor[HTML]{EAEAEA}\textbf{\$36.443.570,64} \\ \hline
\end{tabular}
\caption{Cálculo del Valor Actual Neto (VAN).}
\label{tab:van}
\end{table}

% --- NUEVA SUBSECCIÓN DE RIESGOS ---
\subsection{Análisis de Resiliencia Económica y Escalabilidad}



\subsubsection*{Consistencia de la inversión} 


\section{Conclusión de factibilidad}
